\documentclass[lang=cn,newtx,11pt,scheme=chinese]{elegantbook}

\title{Probability Theory}
\subtitle{Latex Notes}

\author{YuChen Huang}
\institute{Self Study}
\date{\today}
\version{0.1}
\bioinfo{联系}{0874103232qq.com@gmail.com}



\extrainfo{\textcolor{red}{\bfseries 注意：本模板自2023年1月1日起停止维护。然而，鉴于用户群体庞大，自2026年起恢复维护并重新发布。}}

\setcounter{tocdepth}{3}


 \cover{"c:/Users/Loz20/Desktop/Type Notes/Latex/Latex Notes/Probability Theory/qwq.png"}

% 本文档命令
\usepackage{array}
\newcommand{\ccr}[1]{\makecell{{\color{#1}\rule{1cm}{1cm}}}}

% 修改标题页的橙色带
\definecolor{customcolor}{RGB}{32,178,170}
\colorlet{coverlinecolor}{customcolor}
\usepackage{cprotect}

\begin{document}

\maketitle
\frontmatter

\tableofcontents

\mainmatter
\chapter{Preface}
This note aim to be overall cover through of the textbook of $\textbf{Rick Durrett's "Probability: Theory and Examples"}$. As well as a practice run of the $\textbf{LaTex}$ under a new template.

A pre-requisite of this note is that the reader should have a basic knowledge of measure theory, probability theory. Which I highly recommend the reader to read over MIRA by $\textbf{Sheldon Axler}$ and intro to probability by $\textbf{Rick Durrett}$.

To whoever who used my note and endedup scoring higher than me in the exam, I would like to say fk you. 

August 2025

UCD

\chapter{测度论}


\section{Lebesgue 测度}
\subsection{Lebesgue 测度的定义}
Lebesgue 测度是现代测度论的核心概念之一，它扩展了传统的长度、面积和体积的概念，使得我们能够对更广泛的集合进行测量。Lebesgue 测度的定义基于外测度的概念，通过对集合进行覆盖和取 infimum 来定义其测度。
\begin{definition}[Lebesgue 外测度]
    Lebesgue measure is a measure on $(\mathbb{R}, \mathcal{B})$, that assigns to each Borel set its outer measure, which is defined as the infimum of the lengths of open intervals that cover the set.
\end{definition}

\begin{lemma}{Approximation of Borel sets from above by closed sets}
    For any Borel set $B \subseteq \mathbb{R}$ and any $\epsilon > 0$, there exists a closed set $F \subseteq B$ such that $\mu(B \setminus F) < \epsilon$. 
    
\end{lemma}
这里开始正文。
\chapter{大数定理}
z
\chapter{中心化定理}
z
\chapter{Martingales}
z
\chapter{马科夫链}
zz
\chapter{Browian Motion}
z
\chapter{Random Walk 应用}
z
\chapter{多维Browian 运动}

\chapter{补充：测度论}

\end{document}